% Exercises for Chapter 1 — What Is Reinforcement Learning?

\begin{enumerate}
  \item \textbf{Conceptual.} Explain, in your own words, why supervised
  learning is not a natural fit for teaching a robot to pour coffee.
  What specifically is missing from a supervised-learning dataset that RL
  provides instead?

  \item \textbf{Conceptual.} Give an example (not from this chapter) of a
  real-world problem that is naturally episodic, and one that is
  naturally continuing. Justify each choice.

  \item \textbf{State vs.\ observation.} For a self-driving car, describe
  something that plausibly belongs to the true \emph{state} of the world
  but would \emph{not} be directly available in the car's
  \emph{observation} from its sensors.

  \item \textbf{Numerical.} Using the $1\times4$ Grid World from this
  chapter's worked example, write out the full trajectory (states,
  actions, rewards) if the agent instead takes the actions
  LEFT, RIGHT, RIGHT, RIGHT, RIGHT (recall moving off the strip leaves
  the agent in place). How many total steps does this trajectory take
  to reach the goal, compared to the RIGHT, RIGHT, RIGHT trajectory in
  the chapter?

  \item \textbf{Thought experiment.} Suppose every reward in Grid World
  were changed to exactly $0$, including at the goal. What, if anything,
  could an RL agent learn in this environment? What does this tell you
  about the role reward plays in defining the task itself?

  \item \textbf{Robotics.} Pick one component of the agent-environment
  loop (state, action, reward, or environment) and describe concretely
  what it would correspond to for a warehouse robot picking items off
  shelves. Be specific about units or data type (e.g., "action = a
  3-dimensional continuous vector of end-effector velocity, in m/s").

  \item \textbf{Common mistakes.} A classmate says: "Reward and value are
  basically the same thing, since a bigger reward is always better."
  Explain, using an example, why this statement is misleading.
\end{enumerate}
